%================================================================
\section{Introduction}\label{sec:intro}

A planner wants to know what a campaign would do to preventive behavior in a
village if it seeded the \KSeeded{} best-connected people, compared with no
campaign at all. Seeding the people with the most contacts is a common proposal
for network interventions, and field work has studied both its promise and its
limits \citep{banerjee2013,kim2015}. Before the campaign, the planner can see the
village as it is, and she may be able to run one randomized pilot. She cannot see
the policy itself. A common way to answer the question is to build a simulator of
the village, calibrate it to observable moments, and run the policy inside the
simulator \citep{kydland1996}. The answer is then one number: the policy effect
at the calibrated parameter point, perhaps with a Monte Carlo error bar.

That number can mislead for a reason that has little to do with the simulator's
quality. Calibration chooses one parameter point, but the targets rarely pin one
down. Many points can match the same moments about equally well, a fact that
simulation modellers call equifinality \citep{beven2006} and that economists have
long raised about calibration \citep{hansen1996}. If those points imply different
policy effects, the single number carries an extra, unstated assumption: that the
chosen point is the world. History matching takes the opposite stance. It rules
out the parameter points that the targets reject and keeps the rest, the
not-ruled-out-yet (NROY) set \citep{craig1997,vernon2010,andrianakis2015}. The
natural report is then the \emph{image} of the policy effect over that set. This
is the reporting logic of partial identification \citep{manski2003,tamer2010}:
when the maintained assumptions leave a set of answers, report the set rather
than a convenient point.

The policy question has a second difficulty. People in a village influence each
other, so the outcome of one person depends on who else is seeded. The policy
effect is then a contrast of two whole assignment vectors, the one that seeds the
hubs and the one that seeds no one \citep{halloran1995,hudgens2008,manski2013}.
A randomized pilot seeds people at random. It realizes the hub-seeding vector only
with negligible probability, so an estimate of the policy effect from the pilot
needs two further assumptions: an exposure map, under which a person's outcome
depends on the assignment only through a low-dimensional summary, and transport,
under which the outcomes measured on the pilot's people apply to the people who
would hold the same exposures under the policy. An analyst who ignores
interference avoids both assumptions only by making a stronger one.

This paper puts the two difficulties side by side in a setting where the truth is
known. The village is an agent-based model with \NAgents{} people on a fixed
random graph. The targets are planted: they are the model's own moments at a
chosen point $\cstar$, the real village of the exercise. They are what a
researcher could observe, the status quo and one randomized pilot; the policy to
predict is never a target. History matching keeps the points that the targets
cannot rule out, and we report the image of the policy effect over a declared
sample of them. We then hand one panel to several analysts, each with an interval
for the policy effect (or, in one case, for a spillover contrast), and we grade
each interval in two separate ways: whether it contains the truth, and how it sits
against the image. The analysts form a ladder of assumptions. The first assumes
no interference. Two more allow for interference through an exposure map and
transport from the pilot, the second of them without borrowing control outcomes
across a split of the village into hubs and everyone else. They work for two
kinds of researcher: a field experimenter who has one panel from the real village,
and a calibrator who works at one best-fitting point of the model.

The contribution can be stated in three sentences. First, we write the policy
effect under interference as a contrast of two assignment vectors and make
explicit what an estimate of it from a randomized pilot must assume. Second, we
report the image of that effect over the parameter points that observable targets
do not rule out, and we grade one-panel intervals against both the truth and the
image, keeping containment and interval geometry apart. Third, we measure over
repeated panels what an analyst gains and loses by allowing for interference:
less bias, more spread, and intervals that are conservative relative to their
sampling distribution.

% HEADLINE: adopted by the owner, 2026-11-01 (PM wording, checked by verifier V6)
\begin{statusbox}[Headline]
In the one panel at the real village, the analyst who assumes no interference
reports a tight interval that misses the policy effect. At the calibrator's two
best-fitting points, the same analyst's interval misses again and is
overconfident: it lies strictly inside the range of effects that the targets
leave open. Over repeated panels, analysts who allow for interference through an
exposure map and transport from the pilot reduce the absolute bias at every
point, at the price of a larger sampling spread and intervals that are
conservative relative to it. Their intervals still answer only for the world in
which the analyst works: in the one panel at the median-I point, they contain
that point's effect but not the village's. The image contains the effect at the
village and at both of the calibrator's points.
Report the image, and read any one-point interval against it.
\end{statusbox}
% END HEADLINE

Two words in the headline carry precise meanings, and we keep them apart
throughout. An interval \emph{misses} when it does not contain the true value at
the point in question; this is a statement about one realized panel and one
world. An interval is \emph{overconfident}, in the sense of the verdict
F-OVERCONF (Section~\ref{sec:geometry}), when it lies strictly inside the image:
it is narrower than the set of effects that the targets leave open. An interval
can contain the truth and still be overconfident, and it can miss the truth
without being overconfident. We apply the charge of overconfidence only to the
calibrator, who adds the assumption that her point is the world. The field
experimenter works with the real village's panel; we report how her interval sits
against the image, but that comparison is not a test of her estimator.

None of the pieces is new. History matching, partial identification, exposure
mappings and robust standard errors are all standard, and the paper treats each
of them as a well-built tool for the question it was designed to answer. What we
add is their composition around one named policy contrast, in a planted setting
that lets every interval be checked. The numbers are the simulator's own, so the
paper is a methods illustration with no claim about any real campaign.

Section~\ref{sec:estimand} defines the estimand and states what the pilot can and
cannot reveal about it. Section~\ref{sec:model} describes the village and the
targets, and Section~\ref{sec:hm} describes history matching and the image.
Section~\ref{sec:analysts} sets out the ladder of analysts, and
Section~\ref{sec:evaluation} the evaluation. Section~\ref{sec:results} reports the
results, Section~\ref{sec:pi} relates the procedure to partial identification,
Section~\ref{sec:limits} lists what the procedure is not, and
Section~\ref{sec:conclusion} concludes. The appendix holds the model details,
the checks behind each analyst, the sensitivity menu, replication and the
disclosures.
